\documentclass[twoside]{article}

\usepackage[preprint]{aistats2027}
\usepackage[round,authoryear]{natbib}
\usepackage{amsmath,amssymb,amsthm}
\usepackage{booktabs}
\usepackage{balance}
\usepackage{graphicx}
\usepackage{microtype}
\usepackage{url}

\bibliographystyle{apalike}

\newtheorem{theorem}{Theorem}
\newtheorem{proposition}{Proposition}
\newtheorem{lemma}{Lemma}
\newtheorem{corollary}{Corollary}
\newcommand{\E}{\mathbb E}
\newcommand{\Pcal}{\mathcal P}
\newcommand{\TV}{\operatorname{TV}}
\newcommand{\diam}{\operatorname{diam}}

\begin{document}

\twocolumn[
\aistatstitle{Supplement to ``Extrapolation Is an Assumption: Sharp Limits for Pass@k Forecasting''}

\aistatsauthor{}

\aistatsaddress{}
]

\section{Identification Proofs}

\subsection{Observable information}

Let
\[
 B_{s,b}(p)=\binom bs p^s(1-p)^{b-s},
 \qquad s=0,\ldots,b,
\]
be the degree-\(b\) Bernstein basis. If \(S\mid P=p\) is binomial, then
\(\Pr(S=s)=\int B_{s,b}(p)dF(p)\).

\begin{proof}[Proof of Proposition 1]
The \(b+1\) Bernstein polynomials form a basis for \(\Pcal_b\). Hence the
count probabilities determine the integral of every polynomial in
\(\Pcal_b\). Conversely, every linear functional of the count probabilities
is the expectation of a polynomial in their span, so no other linear
functional of \(F\) is directly determined.

For \(k\leq b\), conditional on \(S=s\), the probability that a uniformly
chosen subset of \(k\) observed attempts contains at least one success is
\[
 \widehat h_k(s)=1-\frac{\binom{b-s}{k}}{\binom bk},
\]
with the ratio defined as zero when \(b-s<k\). Conditional exchangeability
gives \(\E_p[\widehat h_k(S)]=1-(1-p)^k\). Averaging over tasks and \(F\)
proves unbiasedness.

For \(k>b\), \(h_k(p)=1-(1-p)^k\) has degree \(k\) and is not in
\(\Pcal_b\). Theorem 1 constructs two laws that agree on every element of
\(\Pcal_b\) but disagree on \(h_k\), proving nonidentification in general.
\end{proof}

\subsection{The identified interval for fixed moments}

For a feasible moment vector
\(\mu=(\mu_0,\ldots,\mu_b)\), define
\[
 \mathcal F(\mu)=
 \left\{F:\int p^j\,dF(p)=\mu_j,\ j=0,\ldots,b\right\}.
\]
The set is weakly compact because probability measures on \([0,1]\) are
weakly compact and the moment constraints are continuous. The functional
\(F\mapsto\int h_kdF\) is continuous, so its minimum and maximum over
\(\mathcal F(\mu)\) exist. Convexity implies that every value between the
two endpoints is attained.

The lower primal problem and its generalized moment dual are
\begin{align*}
\inf_{F\in\mathcal F(\mu)}\int h_k\,dF
\quad\text{and}\quad
\sup_{q\in\Pcal_b}
\left\{\sum_{j=0}^bq_j\mu_j:q\leq h_k\right\}.
\end{align*}
The upper pair reverses the optimization and the envelope. Feasibility and
compact support give equality of primal and dual values
\citep{deklerk2019,halaseh2026}. A dual envelope need not attain its optimum
when \(\mu\) lies on the boundary of the moment cone. This is why the main
paper uses a supremum and an infimum.

Because \(h_k-q\) and \(q-h_k\) are univariate polynomials, their
nonnegativity on \([0,1]\) can also be represented exactly by the
Markov--Lukacs certificate. If a polynomial \(r\) has even degree \(2d\),
\[
 r(p)=\sigma_0(p)+p(1-p)\sigma_1(p),
\]
where \(\sigma_0,\sigma_1\) are sums of squares of degrees at most \(2d\) and
\(2d-2\). If its degree is odd \(2d+1\),
\[
 r(p)=p\sigma_0(p)+(1-p)\sigma_1(p),
\]
where both sums of squares have degree at most \(2d\). Gram-matrix
representations turn the endpoint problems into finite semidefinite programs.
Their size grows with \(k\), so our largest-horizon computations instead use
linear programs on Chebyshev grids and monotone outer projection.

\subsection{Sharp worst-case diameter}

\begin{proof}[Proof of Theorem 1]
Set \(X=1-P\). Matching the first \(b\) moments of \(P\) is equivalent to
matching those of \(X\). Let \(F,G\) be any two such laws and let
\(q\in\Pcal_b\). Since \(\int q\,d(F-G)=0\),
\begin{align*}
\left|\theta_k(F)-\theta_k(G)\right|
&=\left|\int x^k\,d(F-G)(x)\right|\\
&=\left|\int(x^k-q(x))\,d(F-G)(x)\right|\\
&\leq 2\|x^k-q\|_\infty.
\end{align*}
Taking the infimum over \(q\) proves \(D_{b,k}\leq2E_{b,k}\).

For the reverse inequality, let \(C=C[0,1]\) with the uniform norm, and
consider the quotient Banach space \(C/\Pcal_b\). The norm of the coset
\([x^k]\) is its distance to \(\Pcal_b\), namely \(E_{b,k}\). By
Hahn--Banach, there is a continuous linear functional \(\Lambda\) on \(C\)
with operator norm one such that
\[
 \Lambda(q)=0\quad(q\in\Pcal_b),\qquad
 \Lambda(x^k)=E_{b,k}.
\]
The Riesz representation theorem gives a finite signed measure \(\nu\) with
total variation one and \(\Lambda(f)=\int f\,d\nu\). Since the constant
function lies in \(\Pcal_b\), \(\nu([0,1])=0\). In the Jordan decomposition
\(\nu=\nu^+-\nu^-\), both parts therefore have mass \(1/2\).

Define probability measures \(F^\star=2\nu^+\) and
\(G^\star=2\nu^-\). They match every polynomial of degree at most \(b\), and
\[
 \int x^k\,d(F^\star-G^\star)=2E_{b,k}.
\]
Their count laws match by the Bernstein representation. Their pass@\(k\)
values have the same absolute gap because the constant part of
\(1-x^k\) cancels. Thus \(D_{b,k}\geq2E_{b,k}\).
\end{proof}

\subsection{Newman--Rivlin mapping}

\begin{proof}[Proof of Corollary 1]
Let \(E_n^{[-1,1]}(f)\) denote best degree-\(n\) uniform approximation on
\([-1,1]\). If \(q\in\Pcal_b\), then \(q(z^2)\) is an even degree-\(2b\)
polynomial and
\[
 \sup_{z\in[-1,1]}|z^{2k}-q(z^2)|
 =\sup_{x\in[0,1]}|x^k-q(x)|.
\]
Conversely, symmetrizing any degree-\(2b\) approximant to the even function
\(z^{2k}\) cannot increase its error. The resulting even polynomial is
\(q(z^2)\) for a \(q\in\Pcal_b\). Hence
\[
 E_{b,k}=E_{2b}^{[-1,1]}(z^{2k}).
\]

Theorem 3 of \citet{newman1976} states that
\[
 \frac{1}{4e}P_{n,m}\leq E_m^{[-1,1]}(z^n)\leq P_{n,m},
\]
where \(P_{n,m}\) is the probability that the absolute difference between
heads and tails in \(n\) fair tosses exceeds \(m\). With \(n=2k,m=2b\),
that event is \(|H-k|>b\) for
\(H\sim\operatorname{Binomial}(2k,1/2)\). Multiplication by two and Theorem 1
give the displayed bounds.

For large \(k\), the standardized threshold is
\(b/\sqrt{k/2}=\sqrt{2}b/\sqrt{k}\). The central limit theorem therefore
places the transition at \(b/\sqrt{k}\) of constant order. If
\(b/\sqrt{k}\to\infty\), the tail and diameter vanish. If that ratio remains
bounded, the lower bound stays away from zero along nondegenerate limits.
\end{proof}

\subsection{Declared support restrictions}

\begin{proof}[Proof of Corollary 2]
Under \(P\in[p_0,1]\), set \(X=1-P\in[0,a]\) with \(a=1-p_0\). If \(a=0\),
the target is constant and the claim is immediate. Otherwise write \(X=aY\).
The map from laws of \(X\) to laws of \(Y\in[0,1]\) is bijective. Matching
degree-\(b\) moments is preserved under this scaling, and
\[
 \left|\E_FX^k-\E_GX^k\right|
 =a^k\left|\E_FY^k-\E_GY^k\right|.
\]
Taking the supremum over matching laws and applying Theorem 1 gives
\(D_{b,k}(p_0)=a^kD_{b,k}\).
\end{proof}

\section{Confidence-Set Lower Bounds}

\subsection{Exactly matching count laws}

\begin{proof}[First part of Theorem 2]
Let \(A_F=\{\theta_k(F)\in C\}\) and
\(A_G=\{\theta_k(G)\in C\}\). Under a shared data law,
\[
 \Pr(A_F\cap A_G)
 \geq1-\Pr(A_F^c)-\Pr(A_G^c)\geq1-2\alpha.
\]
On the intersection, \(C\) contains two points separated by \(\Delta\), so
its diameter is at least \(\Delta\). The expected-length claim follows from
\(\E[\diam(C)]\geq\Delta\Pr\{\diam(C)\geq\Delta\}\).
\end{proof}

\subsection{Nearby finite-task experiments}

\begin{lemma}[Affinity bound]
For probability mass functions \(q,r\), define
\(\rho(q,r)=\sum_s\sqrt{q_sr_s}\). Then
\[
 \TV(q,r)\leq\sqrt{1-\rho(q,r)^2}.
\]
For \(M\) iid observations,
\(\rho(q^M,r^M)=\rho(q,r)^M\).
\end{lemma}

\begin{proof}
Cauchy--Schwarz gives
\begin{align*}
\TV(q,r)
&=\frac12\sum_s|\sqrt{q_s}-\sqrt{r_s}|
                    (\sqrt{q_s}+\sqrt{r_s})\\
&\leq\frac12
\sqrt{\sum_s(\sqrt{q_s}-\sqrt{r_s})^2}
\sqrt{\sum_s(\sqrt{q_s}+\sqrt{r_s})^2}\\
&=\sqrt{(1-\rho(q,r))(1+\rho(q,r))}.
\end{align*}
The product identity follows by factorizing the sum over product outcomes.
\end{proof}

\begin{proof}[Finite-task part of Theorem 2]
Let \(P_F,P_G\) denote the two full data laws. Coverage and total variation
give
\begin{align*}
 P_F(A_G)
 &\geq P_G(A_G)-\TV(P_F,P_G)\\
 &\geq1-\alpha-\TV(P_F,P_G).
\end{align*}
Combining this with \(P_F(A_F)\geq1-\alpha\) yields
\[
 P_F(A_F\cap A_G)\geq1-2\alpha-\TV(P_F,P_G).
\]
The preceding lemma bounds the final term by
\(\sqrt{1-\rho(q_F,q_G)^{2M}}\). On the intersection the interval spans the
target gap.
\end{proof}

\subsection{Convex modulus program and repair}

On support points \(p_1,\ldots,p_J\), let \(w,z\) be two vectors in the
probability simplex. Define the \((b+1)\times J\) matrix
\[
 A_{sj}=\binom bs p_j^s(1-p_j)^{b-s}
\]
and \(t_j=h_k(p_j)\). For a desired product-TV budget \(v\), the
grid-restricted program is
\begin{align*}
\text{maximize}\quad &t^\top(w-z)\\
\text{subject to}\quad
&w,z\geq0,\quad\mathbf1^\top w=\mathbf1^\top z=1,\\
&\sum_{s=0}^b\sqrt{(Aw)_s(Az)_s}
\geq(1-v^2)^{1/(2M)}.
\end{align*}
The affinity is concave in \((w,z)\), so its superlevel set is convex. The
objective is linear. Any feasible solution is an actual pair of mixing laws,
and therefore a valid lower-bound witness for the unrestricted modulus.

Numerical solvers can return weights with tiny negative entries or affinity
slightly below the requested value. We first clip negative weights and
renormalize. Let \(m=(w+z)/2\) and \(d=(w-z)/2\). We then replace the pair by
\[
 w_\lambda=m+\lambda d,\qquad z_\lambda=m-\lambda d,
 \qquad0\leq\lambda\leq1.
\]
At \(\lambda=0\), the two laws coincide and have affinity one. We use
bisection to select the largest feasible \(\lambda\), followed by a relative
\(10^{-10}\) shrinkage. Every reported target gap is recomputed from this
repaired pair. Thus solver optimality labels do not affect validity of the
lower bounds.

\begin{table*}[t]
\caption{Finite-task modulus witnesses for \(M=128,b=78,\alpha=0.05\).
All displayed product-TV bounds are at most the requested \(v\).}
\label{tab:modulus-full}
\centering
\small
\begin{tabular}{rrrrrr}
\toprule
\(k\) & \(v\) & target gap & span probability & expected-length lower & solver label\\
\midrule
100 & .05 & .00531 & .850 & .00451 & optimal-inaccurate\\
100 & .10 & .01063 & .800 & .00850 & optimal\\
100 & .20 & .02139 & .700 & .01497 & optimal\\
100 & .30 & .03246 & .600 & .01948 & optimal\\
100 & .40 & .04407 & .500 & .02203 & optimal\\
\midrule
1,000 & .05 & .28233 & .850 & .23998 & optimal-inaccurate\\
1,000 & .10 & .33187 & .800 & .26549 & optimal-inaccurate\\
1,000 & .20 & .39457 & .700 & .27620 & optimal-inaccurate\\
1,000 & .30 & .44060 & .600 & .26436 & optimal-inaccurate\\
1,000 & .40 & .47946 & .500 & .23973 & optimal-inaccurate\\
\midrule
10,000 & .05 & .80759 & .850 & .68645 & optimal-inaccurate\\
10,000 & .10 & .83681 & .800 & .66945 & optimal-inaccurate\\
10,000 & .20 & .86959 & .700 & .60871 & optimal-inaccurate\\
10,000 & .30 & .89134 & .600 & .53480 & optimal-inaccurate\\
10,000 & .40 & .90786 & .500 & .45393 & optimal-inaccurate\\
\bottomrule
\end{tabular}
\end{table*}

\subsection{Adaptive policies}

\begin{proof}[Proof of Proposition 2]
Condition on the policy's external random seed, so its decisions are
deterministic functions of the observed transcript. Consider any finite
transcript compatible with the policy. For task \(i\), suppose it contains
\(t_i\leq b\) observations with \(s_i\) successes. Integrating its latent
probability gives the history probability
\[
 \int p^{s_i}(1-p)^{t_i-s_i}\,dF(p).
\]
The integrand has degree \(t_i\leq b\), so this value is unchanged when
\(F\) is replaced by a law \(G\) with matching first \(b\) moments. Latent
task probabilities are independent, hence the probability of the collection
of task histories is the product of these integrals. Policy decisions add
only indicators that the transcript follows the deterministic rule.
Therefore every complete transcript has the same probability under \(F\)
and \(G\). Averaging over the external random seed proves the claim for a
randomized policy.
\end{proof}

\section{Honest Interval Constructions}

\subsection{Global Bernstein interval}

\begin{theorem}
Fix arbitrary \(p_1,\ldots,p_M\in[0,1]\), and independently draw
\(S_i\sim\operatorname{Binomial}(b,p_i)\). If
\(\|q-h_k\|_\infty\leq\epsilon\), where
\[
 q(p)=\sum_{s=0}^ba_sB_{s,b}(p),
\]
and \(R=\max_sa_s-\min_sa_s\), then the global polynomial interval in the main
paper covers
\[
 \theta_{M,k}=\frac1M\sum_{i=1}^Mh_k(p_i)
\]
with probability at least \(1-\alpha\). If instead \(P_i\) are iid from \(F\),
the same interval also covers population \(\theta_k(F)\) with that probability.
\end{theorem}

\begin{proof}
Let \(\bar a=M^{-1}\sum_i a_{S_i}\). The Bernstein identity gives
\[
 \E[\bar a]=\frac1M\sum_iq(p_i).
\]
The deterministic approximation error therefore satisfies
\[
 |\E[\bar a]-\theta_{M,k}|\leq\epsilon.
\]
Each \(a_{S_i}\) lies in an interval of length \(R\). Hoeffding's inequality
for independent, not necessarily identically distributed variables yields
\[
 \Pr\{|\bar a-\E\bar a|>Rc_M\}\leq
 2\exp(-2Mc_M^2)=\alpha.
\]
The triangle inequality gives coverage before clipping, and intersecting with
\([0,1]\) cannot remove the target.

For iid tasks, the mixture counts \(S_i\) are iid and
\(\E[a_{S_i}]=\int q(p)dF(p)\). Repeating the same approximation and
concentration argument gives population coverage.
\end{proof}

\paragraph{Design program.}
On evaluation points \(u_\ell\), the implemented linear program minimizes
\[
 \epsilon+c_M(a_{\max}-a_{\min})
\]
over \(a_0,\ldots,a_b,\epsilon,a_{\min},a_{\max}\), subject to
\[
 |q(u_\ell)-h_k(u_\ell)|\leq\epsilon,\qquad
 a_{\min}\leq a_s\leq a_{\max}.
\]
After solving, we certify the continuous error. The Bernstein derivative
identity gives
\[
q'(p)=b\sum_{s=0}^{b-1}(a_{s+1}-a_s)B_{s,b-1}(p),
\]
so \(|q'(p)|\leq b\Delta_a\), where
\(\Delta_a=\max_s|a_{s+1}-a_s|\). Also
\(h_k'(p)=k(1-p)^{k-1}\), which decreases in \(p\). On adjacent grid points
\(u_\ell,u_{\ell+1}\), every \(p\) is within half the interval width of an
endpoint. Hence the continuous error on that interval is at most
\begin{align*}
\max\{&
|q(u_\ell)-h_k(u_\ell)|,\,
|q(u_{\ell+1})-h_k(u_{\ell+1})|\}\\
&+\frac{u_{\ell+1}-u_\ell}{2}
\left[k(1-u_\ell)^{k-1}+b\Delta_a\right].
\end{align*}
We use the maximum of these bounds plus \(10^{-10}\) as \(\epsilon\).

\paragraph{Asymptotic sharpness.}
In the ideal continuous design problem, every feasible polynomial has
\(\epsilon\geq E_{b,k}\). Conversely, for any \(\eta>0\), choose a polynomial
with error at most \(E_{b,k}+\eta\). Its finite Bernstein coefficient range is
multiplied by \(c_M\to0\). Therefore the optimized half-width converges to
\(E_{b,k}\), and the full worst-case width converges to
\(2E_{b,k}=D_{b,k}\).

\subsection{Count-CDF projection interval}

For thresholds \(s=0,\ldots,b-1\), write
\[
 g_s(p)=\Pr\{\operatorname{Binomial}(b,p)\leq s\}.
\]
This function is decreasing in \(p\). The pass target \(h_k(p)\) is
increasing.

\begin{theorem}
For iid tasks from \(F\), use radius
\[
 \delta_{\rm iid}=\sqrt{\frac{\log(2/\alpha)}{2M}}.
\]
For fixed \(p_1,\ldots,p_M\) with independent counts, use
\[
 \delta_{\rm fix}=\sqrt{\frac{\log(2b/\alpha)}{2M}}.
\]
The endpoint-envelope linear program in the main paper gives a
\((1-\alpha)\) outer confidence interval for \(\theta_k(F)\) in the iid case,
and for \(\theta_{M,k}\) in the fixed case.
\end{theorem}

\begin{proof}
In the iid case, the counts \(S_i\) are iid draws from the mixture count CDF
\(G_s=\int g_s(p)dF(p)\). The DKW inequality
\citep{massart1990} gives
\[
 \Pr\left\{\max_s|\widehat G_s-G_s|>\delta_{\rm iid}\right\}
\leq\alpha.
\]

In the fixed case,
\(\mathbf1\{S_i\leq s\}\) are independent variables in \([0,1]\), with
average expectation
\[
 G_{M,s}=\frac1M\sum_i g_s(p_i).
\]
For each threshold, Hoeffding's inequality gives a two-sided error probability
at most \(2\exp(-2M\delta^2)\). A union bound over the \(b\) nontrivial
thresholds gives the stated radius.

It remains to prove projection validity on either simultaneous-band event.
Let \(F_\star\) be \(F\) in the iid setting and the empirical law
\(M^{-1}\sum_i\delta_{p_i}\) in the fixed setting. On a partition
\([e_j,e_{j+1}]\), set \(w_j=F_\star([e_j,e_{j+1}])\), with a fixed convention
at shared endpoints. Monotonicity gives
\[
 \sum_jw_jg_s(e_{j+1})
 \leq\int g_s\,dF_\star
 \leq\sum_jw_jg_s(e_j).
\]
Combining these inequalities with the confidence band shows that the true
bin masses are feasible for the linear constraints. A second use of
monotonicity gives
\[
 \sum_jw_jh_k(e_j)
 \leq\int h_k\,dF_\star
 \leq\sum_jw_jh_k(e_{j+1}).
\]
The minimized lower envelope can only be smaller than the left side, and the
maximized upper envelope can only be larger than the right side. The projected
interval therefore contains the target on the simultaneous-band event.
\end{proof}

\paragraph{Numerical conditioning.}
To compute the upper endpoint, maximizing \(h_k\) is implemented as minimizing
the failure probability \((1-p)^k\), then subtracting the optimum from one.
This avoids an objective numerically indistinguishable from the constant
one at large \(k\). Each endpoint first uses HiGHS, then retries its dual
simplex and interior-point methods at a relaxed \(10^{-8}\) feasibility
tolerance. Every reported run completed successfully after this
reparameterization.

\section{Numerical Identification Results}

\begin{table}[t]
\caption{Worst-case population identification. Grid-based approximation and
moment-pair programs agree at the displayed precision.}
\label{tab:diameters}
\centering
\small
\begin{tabular}{rrrr}
\toprule
\(b\) & \(k\) & approx.\ \(D_{b,k}\) & max moment resid.\\
\midrule
4  & 128    & .6677 & \(2.1\times10^{-15}\)\\
8  & 1,000  & .7833 & \(5.7\times10^{-12}\)\\
16 & 1,000  & .5086 & \(5.1\times10^{-15}\)\\
32 & 10,000 & .7066 & \(4.3\times10^{-12}\)\\
64 & 10,000 & .3886 & \(1.6\times10^{-14}\)\\
78 & 1,000  & .000489 & \(8.2\times10^{-11}\)\\
78 & 10,000 & .28495 & \(7.3\times10^{-13}\)\\
\bottomrule
\end{tabular}
\end{table}

The \(b=78,k=100\) error is below double-precision resolution. We do not use
the numerical zero as a theorem. Corollary 1 gives the rigorous bound
\[
 3.58\times10^{-33}\leq D_{78,100}
 \leq3.89\times10^{-32}.
\]

\begin{table*}[t]
\caption{Atomic laws for the \(b=8,k=1000\) ambiguity example. Weights shown
here are rounded. The artifact stores full precision.}
\label{tab:atoms}
\centering
\small
\begin{tabular}{rrrr}
\toprule
\multicolumn{2}{c}{High pass@1000 law} &
\multicolumn{2}{c}{Low pass@1000 law}\\
\(p\) & weight & \(p\) & weight\\
\midrule
.002934 & .940617 & .000000 & .833112\\
.149794 & .038257 & .041906 & .134095\\
.501963 & .011496 & .311564 & .018474\\
.854108 & .006747 & .692430 & .008329\\
1.000000 & .002882 & .962090 & .005991\\
\bottomrule
\end{tabular}
\end{table*}

At full precision, the two induced count probability vectors differ by at
most \(1.7\times10^{-12}\). Their target values are \(0.9501871171\) and
\(0.1668880229\). Their stored degree-\(b\) Chebyshev moment residual is
\(5.7\times10^{-12}\). For this ambiguity example, the numerical primal and
dual diameters differ by less than \(10^{-11}\).

\section{Experiment Protocols and Additional Results}

\subsection{Synthetic interval audit}

Each repetition first samples 128 latent task probabilities from the specified
discrete or Beta law. It then samples one binomial count per task. Coverage is
evaluated against the realized finite-task target
\[
 \theta_{M,k}=M^{-1}\sum_{i=1}^M[1-(1-p_i)^k].
\]
This makes the comparison conditional on the sampled benchmark and matches
the fixed-task guarantee of both honest intervals.

The exact pair is the law in Table~\ref{tab:atoms}. Finite-pair laws are the
repaired \(v=0.2\) modulus witnesses for \(b=78,k\in\{1000,10000\}\).
The benign laws are \(\operatorname{Beta}(1,1)\) and
\(\operatorname{Beta}(0.35,5)\), discretized by 20,001 equal-mass midpoint
quantiles for simulation.

For each taskwise Bayesian interval, the posterior is
\[
 p_i\mid S_i\sim\operatorname{Beta}(S_i+1,b-S_i+1).
\]
We use 2,000 independent joint posterior draws, compute the task-average
pass@\(k\) for each draw, and take its 2.5\% and 97.5\% quantiles. The
posterior mean is computed analytically. Population Beta point estimates use
a ten-start Beta-binomial maximum-likelihood fit in log parameters. We also
report the direct plug-in and the clipped binomial smoothed Good--Toulmin
point estimates.

\begin{table*}[t]
\caption{Synthetic 95\% interval audit over 400 repetitions per scenario.
``Task B.'' is taskwise Beta, ``Projection'' is fixed-task count-CDF
projection, and ``Global'' is the Bernstein interval.}
\label{tab:synthetic-full}
\centering
\small
\begin{tabular}{llrrrr}
\toprule
Scenario & Method & coverage & mean width & MAE & repetitions\\
\midrule
Uniform, \(b=78,k=1000\) & Task B. & .9625 & .0059 & .0011 & 400\\
 & Projection & 1.0000 & .1901 & .0941 & 400\\
 & Global & 1.0000 & .7523 & .2606 & 400\\
\midrule
Rare Beta, \(b=78,k=1000\) & Task B. & .0000 & .0379 & .1233 & 400\\
 & Projection & 1.0000 & .5470 & .1224 & 400\\
 & Global & 1.0000 & .8056 & .1637 & 400\\
\midrule
Exact low, \(b=8,k=1000\) & Task B. & .0000 & .0210 & .8235 & 400\\
 & Projection & 1.0000 & 1.0000 & .3318 & 400\\
 & Global & 1.0000 & 1.0000 & .3318 & 400\\
Exact high, \(b=8,k=1000\) & Task B. & .0000 & .0210 & .0415 & 400\\
 & Projection & 1.0000 & .9999 & .4502 & 400\\
 & Global & 1.0000 & 1.0000 & .4503 & 400\\
\midrule
Finite first, \(b=78,k=1000\) & Task B. & .0000 & .0549 & .1384 & 400\\
 & Projection & 1.0000 & .9472 & .2792 & 400\\
 & Global & 1.0000 & .9531 & .2670 & 400\\
Finite second, \(b=78,k=1000\) & Task B. & .0000 & .0551 & .5350 & 400\\
 & Projection & 1.0000 & .9490 & .1165 & 400\\
 & Global & 1.0000 & .9540 & .1285 & 400\\
\midrule
Finite first, \(b=78,k=10000\) & Task B. & .0000 & .0198 & .0204 & 400\\
 & Projection & 1.0000 & 1.0000 & .4722 & 400\\
 & Global & 1.0000 & 1.0000 & .4722 & 400\\
Finite second, \(b=78,k=10000\) & Task B. & .0000 & .0198 & .8900 & 400\\
 & Projection & 1.0000 & 1.0000 & .3974 & 400\\
 & Global & 1.0000 & 1.0000 & .3974 & 400\\
\bottomrule
\end{tabular}
\end{table*}

The honest intervals' empirical coverage of one is consistent with their
guarantees but should not be interpreted as exact calibration. Their
concentration radii are conservative at \(M=128\). The comparison of widths
is the relevant positive result: the projection interval adapts substantially
under Uniform and rare Beta laws, while retaining the required width on hard
pairs.

\subsection{Public response-bank audit}

\paragraph{Data provenance.}
We use the correctness indicators from all 22 configurations of
\texttt{ScalingIntelligence/monkey\_business}, the response-bank release of
\citet{brown2024}. The Hugging Face dataset revision observed during the
audit was
\texttt{a9f8f73bcd6948a57ed922cba4e48062ef95f553}. Its dataset card declares
the MIT license. We fetch only the \texttt{is\_corrects} and original task
index columns from the hosted Parquet shards, sort by original task index, and
store Boolean arrays plus source URLs. The 22 panels contain 127--140 tasks
and 10,000 attempts per task.

\paragraph{References and resampling.}
For a response bank with \(N=10000\) attempts and \(c_i\) successful attempts
on task \(i\), the finite-bank pass reference is
\[
 \widehat\theta_k=
 \frac1M\sum_{i=1}^M\left[
 1-\frac{\binom{N-c_i}{k}}{\binom Nk}\right].
\]
This is the expected without-replacement pass@\(k\) of a uniformly sampled
subset from the bank. For each of 50 seeds, we independently shuffle each
task's responses. Equal allocation uses the first
\(\lfloor B/M\rfloor\) responses per task. Dynamic allocation implements
Algorithm 1 of \citet{kazdan2025}: it initializes each task, then allocates
attempts according to the prescribed function of failures and successes. The
nominal budgets are 1,000, 3,000, and 10,000. Targets are 100, 1,000, and
10,000.

\paragraph{Population Beta fit.}
Given possibly unequal \(b_i\) and counts \(s_i\), the likelihood is
\[
 \prod_i
 \binom{b_i}{s_i}
 \frac{\mathrm B(s_i+\alpha,b_i-s_i+\beta)}
      {\mathrm B(\alpha,\beta)}.
\]
The combinatorial factors do not affect optimization. We optimize
\(\log\alpha,\log\beta\) with analytic gradients from ten starting points,
using bounds \([-14,18]\) on both log parameters. Boundary fits are retained
and flagged. Forecasts use
\[
 1-\frac{\mathrm B(\alpha,\beta+k)}
          {\mathrm B(\alpha,\beta)}.
\]
All 59,400 forecast rows used in the main point-estimate audit have finite
objectives. The artifact preserves convergence, gradient, and boundary
diagnostics.

\begin{table}[t]
\caption{Uniform-allocation public point forecasts at total budget 10,000.
Each row contains 1,100 panel-resamples.}
\label{tab:public-full}
\centering
\scriptsize
\setlength{\tabcolsep}{3pt}
\begin{tabular}{llrrr}
\toprule
\(k\) & method & MAE & RMSE & fraction \(>|.1|\)\\
\midrule
100 & Population Beta & .0175 & .0225 & .000\\
100 & SGT & .0161 & .0223 & .000\\
100 & Plug-in & .0396 & .0516 & .055\\
100 & Taskwise Beta & .2901 & .3451 & .729\\
\midrule
1,000 & Population Beta & .0414 & .0558 & .092\\
1,000 & SGT & .2000 & .2764 & .605\\
1,000 & Plug-in & .1481 & .1845 & .511\\
1,000 & Taskwise Beta & .3843 & .4846 & .682\\
\midrule
10,000 & Population Beta & .0697 & .1300 & .170\\
10,000 & SGT & .3529 & .4688 & .713\\
10,000 & Plug-in & .2365 & .2869 & .737\\
10,000 & Taskwise Beta & .3326 & .4489 & .682\\
\bottomrule
\end{tabular}
\end{table}

\begin{table}[t]
\caption{Fixed-task projection audit over 22 panels and 50 resamples.}
\label{tab:projection-public}
\centering
\scriptsize
\setlength{\tabcolsep}{2pt}
\begin{tabular}{rrrrr}
\toprule
\(k\) & ref.\ inclusion & mean width & median width & width \(>.9\)\\
\midrule
100 & 1.000 & .371 & .395 & .000\\
1,000 & 1.000 & .714 & .787 & .405\\
10,000 & 1.000 & .721 & .787 & .405\\
\bottomrule
\end{tabular}
\end{table}

Reference inclusion is not a frequentist coverage estimate. The full bank is
finite, contains the sampled prefix, and estimates rather than equals a latent
target. We report it as a reproducible descriptive check. The synthetic study
provides the controlled coverage experiment.

\subsection{Compute and software}

All optimization and data experiments ran in Python 3.10 on Argonne Swing.
CPU computations were submitted to a GPU-only PBS queue but do not use the
allocated accelerator. Core packages are NumPy, SciPy 1.15.3, pandas,
Matplotlib, CVXPY, Clarabel, PyArrow, and fsspec. Four spawned workers run the
coverage and projection audits with all BLAS thread counts set to one.

The full 400-repetition synthetic audit took 2 minutes 41 seconds on one node.
The 3,300-row public projection audit took 2 minutes 16 seconds. The artifact
uses fixed seeds, writes outputs only after all worker results return, and
contains ten deterministic tests. Tests cover the combinatorial estimator,
dynamic allocation, Beta-binomial recovery, taskwise posterior computation,
smoothed Good--Toulmin interpolation, identification primal-dual agreement,
continuous bias certification, finite-modulus feasibility, and projection
containment.

\section{Additional Literature Boundary}

The closest statistical model predates pass@\(k\). \citet{colwell2004} use the
same incidence binomial mixture and accumulation functional. They recommend
moment interpolation, fit finite mixtures for extrapolation, and caution that
the asymptote is unreliable. \citet{orlitsky2016} derive minimax smoothed
Good--Toulmin estimators in both abundance and Bernoulli-product incidence
models. Their estimable extrapolation factor grows logarithmically with sample
size. Our work addresses a different quantity: the exact finite-\(b\)
identified diameter, finite-task interval lower bounds, and honest pass@\(k\)
projection.

\citet{mao2007} show that unseen-class odds in a zero-truncated Poisson
mixture are discontinuous and that uniformly informative two-sided inference
for total richness is impossible. That target is unbounded and asymptotic.
Our finite-horizon target remains in \([0,1]\), producing a bounded but
nontrivial ambiguity that changes continuously with \(b\) and \(k\).
\citet{daley2013} use rational continuation for sequencing-library complexity,
while \citet{shestopaloff2025} fits parametric accumulation-rate tails and
uses a parametric bootstrap. Both demonstrate that structural continuation
can be practically effective.

\citet{basu2026} give exact pointwise CDF and quantile intervals for a
binomial mixing distribution by test inversion. They explicitly discuss the
finite-moment identification boundary. They do not study the pass@\(k\)
functional, its sharp diameter, or the \(b^2\) horizon.
\citet{kaido2019} provide generic asymptotically calibrated projection
intervals for partially identified moment models. Our count-CDF interval is a
specialized nonasymptotic outer projection whose validity follows from DKW or
finite Hoeffding bands and monotone bin endpoints.

Recent pass@\(k\) papers impose useful structure rather than study
nonparametric identification. \citet{schaeffer2025} explain aggregate power
laws through heavy lower tails. \citet{levi2025} derive the same curve under a
Beta failure model, while \citet{levi2026} connect a regularly varying
difficulty tail to training-dependent inference scaling. \citet{kazdan2025}
propose population Beta-binomial fitting and dynamic allocation.
\citet{tailpass2026} propagate taskwise Beta or Dirichlet posteriors through
richer repeated-sampling profiles. \citet{hariri2026} use taskwise Bayesian
inference to stabilize model evaluation. Our audits target the extrapolative
assumptions in these population and taskwise models, not their within-model
posterior algebra. A separate diagnostic by \citet{zhou2026} shows that a
deterministic intervention can reach some tasks missed by ordinary sampling.
That result concerns the choice of decoding regime, not extrapolation from a
fixed regime's finite count law.

\balance
\bibliography{references}

\end{document}
